{% Q2081

\needspace{12\baselineskip} 
\item \rtask \ponto{\pt} % 
O método \lstinline[style=Python]|capitalize| da classe \textit{String} do Python é utilizado para: 

\begin{answerlist}[label={\texttt{\Alph*}.},leftmargin=*]
    \ti remover todos os espaços de uma \textit{string}.
    
    \ti verificar se todos os caracteres da \textit{string} são numéricos.
    
    \ti procurar uma substring em uma \textit{string} retornando seu índice caso seja encontrada.
    
    \di retornar uma cópia de uma \textit{string} somente com o primeiro caractere em maiúsculo.
    
    \ti retornar uma cópia de uma \textit{string} com todos os caracteres em minúsculo.
\end{answerlist}

    % A alternativa correta é a letra D, que afirma que o método capitalize da classe String do Python é utilizado para retornar uma cópia de uma string somente com o primeiro caractere em maiúsculo. Este método é bastante simples, mas muito útil em diversas situações onde se deseja garantir que uma string comece com uma letra maiúscula, como, por exemplo, no caso de nomes próprios ou início de sentenças.
    % 
    % Vamos entender melhor o que o método capitalize faz: ele converte o primeiro caractere de uma string para letra maiúscula (se for uma letra alfabética) e o restante dos caracteres para minúsculas, independentemente de como estavam anteriormente. Um exemplo de seu uso seria:
    % 
    % texto = "python é uma linguagem."
    % texto_capitalizado = texto.capitalize()
    % print(texto_capitalizado)  # Saída: "Python é uma linguagem."
    % 
    % Essa funcionalidade é exclusiva para o primeiro caractere da string. Se o primeiro caractere não for uma letra, a função não terá efeito sobre ele, mas ainda assim convertendo o restante dos caracteres para minúsculas. Por exemplo:
    % 
    % texto = "123abc"
    % texto_capitalizado = texto.capitalize()
    % print(texto_capitalizado)  # Saída: "123abc"
    % 
    % Com isso, fica claro por que as outras alternativas são incorretas:
    % 
    %     A fala sobre remover espaços, o que não é a função do capitalize, mas sim de métodos como strip, rstrip ou lstrip.
    %     B menciona a verificação de caracteres numéricos, o que seria o propósito de métodos como isdigit ou isnumeric.
    %     C trata da procura de substrings, que é feita com métodos como find ou index.
    %     E descreve a conversão de todos os caracteres para minúsculas, o que seria feito pelo método lower.
    % 
    % Portanto, ao lembrar que o método capitalize altera somente o caso do primeiro caractere da string (para maiúscula) e o restante dos caracteres para minúsculas, é possível identificar facilmente a resposta correta para essa questão.
}
